% Part: first-order-logic
% Chapter: tableaux
% Section: rules-and-proofs

\documentclass[../../../include/open-logic-section]{subfiles}

\begin{document}

\iftag{FOL}
      {\olfileid{fol}{tab}{rul}}
      {\olfileid{pl}{tab}{rul}}

\olsection{قاعدے تے \usetoken{P}{tableau}}

!!^a{tableau}، !!a{sentence} دے کسے !!a{structure} وچ سچ یا جھوٹ
ہون دے ممکنہ طریقیاں دا منظم جائزہ اے۔ تابلو دیاں بنیادی اکائیاں
!!{signed formula}s نیں: !!{sentence}s دے نال صداقتی قدر دی اک
``علامت''، یعنی $\True$ یا~$\False$۔ ایہ علامت لائے ہوئے
!!{formula}s اک (ہیٹھاں ول ودھدے) درخت وچ ترتیب دتے جاندے نیں۔

\begin{defn}
  اک \emph{!!{signed formula}} صداقتی قدر تے !!a{sentence} دی جوڑی اے، یعنی:
  \[
  \sFmla{\True}{!A} \text{ یا } \sFmla{\False}{!A}.
  \]
\end{defn}

معنی دے لحاظ نال اسیں $\sFmla{\True}{!A}$ نوں ``$!A$ سچ ہو
سکدا اے'' تے $\sFmla{\False}{!A}$ نوں ``$!A$ جھوٹ ہو سکدا اے''
پڑھ سکدے آں (کسے !!{structure} وچ)۔

درخت وچ ہر !!{signed formula} یا تاں \emph{مفروضہ} اے (جیہڑے درخت
دے سب توں اُتّے درج ہوندے نیں)، یا اوہدے اُتّے موجود کسے
!!a{signed formula} توں استنباط دے کئی قاعدیاں وچوں اک ورت کے حاصل
کیتا جاندا اے۔ پچھلے !!{formula} دے ہر ممکن !!{main operator} لئی دو
قاعدے نیں: اک جدوں علامت~$\True$ ہووے، تے اک جدوں علامت~$\False$
ہووے۔ کجھ قاعدے درخت نوں شاخاں وچ ونڈدے نیں، تے کجھ شاخ وچ صرف
!!{signed formula}s شامل کردے نیں۔ قاعدہ فوراً پچھلے
!!{signed formula} اُتے ہی لاگو کرنا ضروری نہیں؛ اوہ جڑ توں اوس جگہ تک والی
شاخ دے کسے وی !!{signed formula} اُتے لاگو ہو سکدا اے (تے اکثر
کرنا پیندا اے)۔

شاخ \emph{بند} اُدوں ہوندی اے جدوں اوہدے وچ $\sFmla{\True}{!A}$
تے $\sFmla{\False}{!A}$ دونیں ہون۔ بند !!{tableau} اوہ اے جیہدی ہر
شاخ بند ہووے۔ بدیہی تشریح مطابق، ہر شاخ اک اکٹھے ممکن ہون والی صورت
دسدی اے، پر $\sFmla{\True}{!A}$ تے $\sFmla{\False}{!A}$ اکٹھے
ممکن نہیں۔ ہور لفظاں وچ، جے شاخ بند اے، تاں اوہ امکان جیہڑا اوہ بیان کردی اے،
ردّ ہو چکیا اے۔ خاص طور تے ایس دا مطلب اے کہ بند !!{tableau} اوہ
سارے امکان ردّ کردا اے جنہاں وچ $\sFmla{\True}{!A}$ دے روپ دے
ہر مفروضے نوں سچ تے $\sFmla{\False}{!A}$ دے روپ دے ہر مفروضے نوں
جھوٹ، اکّو ویلے، بنایا جا سکدا ہووے۔

\emph{$!A$ لئی} بند !!{tableau} اوہ بند !!{tableau} اے جیہدی
جڑ~$\sFmla{\False}{!A}$ ہووے۔ جے ایہو جہا بند !!{tableau} موجود
ہووے، تاں $!A$ دے جھوٹ ہون دے سارے امکان ردّ ہو جاندے نیں؛ یعنی
$!A$ ہر !!{structure} وچ سچ ہونا لازمی اے۔

\end{document}
